\documentclass{leveldoc}
\DocUses{foundations,slingshot,first_coil,rf_kick}
\DocFingerprint{4ac8a9e86ec2cbe8}

\begin{document}

\section{What the level asks}\label{sec:asks}

The game's text:
\leveldescription

One ion is launched from $(\val{xZero}, \val{yZero})$ along $+x$ with $T_0 = \val{Tzero}$, in
Newton's mechanics ($c = \infty$), and must enter the detector
$\val{xDetLo} \le x \le \val{xDetHi}$, $\val{yDetLo} \le y \le \val{yDetHi}$. A circular coil of
radius $a = \val{a}$ about $(\val{xC}, \val{yC})$ lies in the board's plane, without current
until the player sets its power supply, placed on the coil's terminal $(\val{xS}, \val{yS})$. The
supply does not set a current but its rate of rise: the coil's strength (\physref{2.2}:
$\kappa = \mu_0I/4\pi$, in field units) is
\begin{equation}\label{eq:ramp}
  \kappa(t) = \dot\kappa\,t,
\end{equation}
with the rate $\dot\kappa$ on a slider from $\val{rateMin}$ to $\val{rateMax}$. A flight ends when
the ion enters the detector, touches the coil's wire (within $\val{wire}$ of its circle), leaves
the world, or reaches $t = \val{tMax}$ (\xsecref{foundations}{sec:board}).

The reference sets $\dot\kappa = \val{rate}$. The ion passes the coil above it, where the rising
current's field, by then strong, bends it down into the detector. The document derives the
rising coil's fields, the induced electric field among them, marches the flight numerically
with an exact conserved quantity as its check, and finds the interval of rates that solve.

\section{Setup}\label{sec:setup}

The level file, field by field (the appendix reproduces it):
\setuptable

The ion has charge $q = \val{q}$, mass $m = \val{m}$, momentum $p_\mathrm{i} = \val{pZero}$ and
speed $v_0 = \val{vZero}$; it starts $\rho_0 = \val{rhoZero}$ from the coil's centre, and its
launch line passes $\val{rhoLine}$ above the centre, outside the coil. The world is
$-1 \le x \le \val{xWorld}$, $-1 \le y \le \val{yWorld}$.

\begin{figure}[htbp]
\centering
\begin{tikzpicture}
\begin{axis}[axis equal image, width=0.92\linewidth, xmin=-1.6, xmax=31.6, ymin=-1.6, ymax=21.6,
  xtick={0,5,...,30}, ytick={0,5,...,20}, tick label style={font=\small},
  xlabel={$x$}, ylabel={$y$}, label style={font=\small}]
  \draw[black!8] (0,0) grid (30,20);
  \draw[black!45, dashed] (-1,-1) rectangle (\rawval{xWorld},\rawval{yWorld});
  \fill[green!50!black, fill opacity=0.18] (\rawval{xDetLo},\rawval{yDetLo}) rectangle (\rawval{xDetHi},\rawval{yDetHi});
  \draw[green!40!black] (\rawval{xDetLo},\rawval{yDetLo}) rectangle (\rawval{xDetHi},\rawval{yDetHi});
  \draw[black!70, thick] (\rawval{xC},\rawval{yC}) circle[radius=\rawval{a}];
  \fill[yellow!70!black] (\rawval{xS},\rawval{yS}) circle[radius=0.3];
  \addplot[thin, black!50] table {pathlo.dat};
  \addplot[thin, black!50] table {pathhi.dat};
  \addplot[thick, blue!70!black] table {pathref.dat};
  \fill (\rawval{xZero},\rawval{yZero}) circle[radius=0.15];
\end{axis}
\end{tikzpicture}
\caption{The coil (black circle), its supply on the terminal (yellow) and the reference flight
(blue). Grey: the flights at the two edges of the solving interval of \S\ref{sec:others}, grazing
the detector's corners.}
\label{fig:board}
\end{figure}

\section{Principles and model}\label{sec:model}

\begin{principles}
  \item The field of a circular loop in its plane, \xeqref{first_coil}{eq:loop-field}, with $K$
    and $E$ by \xeqref{first_coil}{eq:agm}; the vector potential of a steady current,
    $\vect{A} = \kappa\oint \dd\vect{l}'/R$ in field units, with $\vect{B} = \nabla\times\vect{A}$
    (\Jackson{§5.2, §5.5}); the complete elliptic integrals' power series (DLMF 19.5.1--2).
  \item Faraday's law, the third of \xeqref{foundations}{eq:maxwell}: a changing $\vect{B}$
    comes with the electric field $\vect{E} = -\partial\vect{A}/\partial t$.
  \item Newton's equation with the Lorentz force \xeqref{foundations}{eq:eom} ($c = \infty$), and
    the canonical momentum of a charge in a field, $\vect{p} + q\vect{A}$ (\Jackson{§12.1}),
    whose component along a symmetry of the Lagrangian is conserved (\Goldstein{Ch.~2 and 8}).
  \item The board's rules; the classical Runge--Kutta method, cubic Hermite interpolation and
    the error estimate from the step $2h$ (\xsecref{foundations}{sec:board},
    \xsecref{foundations}{sec:numerics}), as in \xsecref{slingshot}{sec:method} and
    \xsecref{first_coil}{sec:method}.
\end{principles}

The coil's field at each moment is that of its present current \eqref{eq:ramp}, and the rising
current adds the induced electric field $-\dot\kappa\,\vect{A}/\kappa$; there is no other field.
The ion starts at $\vect{x}(0) = (\val{xZero}, \val{yZero})$ with $\vect{p}(0) = (p_\mathrm{i}, 0)$,
when the coil carries no current, and moves until the first of the board's events.

\begin{assumption}{Quasi-static fields}\label{as:quasistatic}
The coil's field follows its current at once, everywhere: the field of a steady current scaled by
$\kappa(t)$, and the induced field $-\partial\vect{A}/\partial t$. This leaves out the coil's
retardation and radiation, an error of the order of the light time across the setup against the
time over which the current changes, $\kappa/\dot\kappa$. In this level's world $c$ is infinite and
the light time is zero, so within the game's model (\physref{2.2}) nothing is neglected. The
supply's own electrostatic field and the wire's charges are not modelled.
\end{assumption}

\begin{assumption}{The coil is a thin loop}\label{as:wire}
The current flows on a circle of zero thickness; the wire's radius $\val{wire}$ only ends a
flight, as in \levelref{first_coil} (its assumption \xref{first_coil}{as:wire}). The reference
flight passes no nearer than $\val{wireGap}$ to the wire.
\end{assumption}

The ion's own radiation needs no assumption: with $c = \infty$, Liénard's power
\xeqref{foundations}{eq:larmor} vanishes.

\section{Derivation}\label{sec:derivation}

\subsection{The loop's vector potential}\label{sec:potential}

Put the loop's centre at the origin and take a point in its plane at distance $\rho$. By symmetry
$\vect{A}$ there points along $\uvec{\phi}$; the wire element at angle $t$ from the point's
direction, $\dd\vect{l}' = a\,\dd t\,(-\sin t, \cos t)$, contributes along $\uvec{\phi}$ its
component $a\cos t\,\dd t$, at distance $R = \sqrt{\rho^2 + a^2 - 2a\rho\cos t}$:
\begin{equation}\label{eq:a-integral}
  A_\phi(\rho) = \kappa a \int_0^{2\pi} \frac{\cos t\,\dd t}{\sqrt{\rho^2 + a^2 - 2a\rho\cos t}}.
\end{equation}
With $t = \pi - 2\theta$, $\cos t = 2\sin^2\theta - 1$ and
$\rho^2 + a^2 - 2a\rho\cos t = \beta^2(1 - m\sin^2\theta)$, where $\beta = a + \rho$ and
$m = 4a\rho/\beta^2$ as in \xeqref{first_coil}{eq:loop-field}; the integrand is even in $\theta$,
so $A_\phi = (4\kappa a/\beta)\int_0^{\pi/2}(2\sin^2\theta - 1)\,\dd\theta/\sqrt{1 - m\sin^2\theta}$.
Writing $2\sin^2\theta - 1 = (2/m - 1) - (2/m)(1 - m\sin^2\theta)$ turns the integral into
$(2/m - 1)K - (2/m)E = (2/m)\,h(m)$, so
\begin{equation}\label{eq:vector-potential}
  A_\phi(\rho) = \frac{2\kappa\beta\,h(m)}{\rho},
  \qquad h(m) = \left(1 - \frac{m}{2}\right)K(m) - E(m),
\end{equation}
\Jackson{§5.5}'s result in the plane of the loop. For small $m$ (near the centre or far away),
$h$ is a difference of nearly equal terms; with the series
$K = \frac{\pi}{2}\sum c_nm^n$, $E = \frac{\pi}{2}\sum c_nm^n/(1 - 2n)$,
$c_n = \big((2n)!/(4^nn!^2)\big)^2$, the terms in $m^0$ and $m^1$ cancel exactly and
\begin{equation}\label{eq:h-series}
  h(m) = \frac{\pi}{2}\sum_{n\ge2} m^n\left(\frac{2n\,c_n}{2n-1} - \frac{c_{n-1}}{2}\right)
       = \frac{\pi m^2}{32} + \frac{3\pi m^3}{128} + \dots,
\end{equation}
used for $m < 0.05$ (where it and the direct form agree to $\val{seriesDiff}$). Two checks: the
trapezoidal rule with 400 nodes on \eqref{eq:a-integral} (spectrally accurate for a smooth
periodic integrand) agrees with \eqref{eq:vector-potential} to $\val{aIntegral}$ relative, and
$B_z = \rho^{-1}\,\dd(\rho A_\phi)/\dd\rho$ by central differences returns
\xeqref{first_coil}{eq:loop-field} to $\val{aCurl}$. Per unit $\kappa$:
\begin{center}\small\val{atab}\end{center}
Outside the loop $B_z$ has the opposite sign to inside and grows in size towards the wire.

\subsection{The rising coil's fields}\label{sec:fields}

With \eqref{eq:ramp}, writing $b(\rho)$ and $a_\phi(\rho)$ for \xeqref{first_coil}{eq:loop-field}
and \eqref{eq:vector-potential} per unit $\kappa$,
\begin{equation}\label{eq:fields}
  B_z = \dot\kappa\,t\;b(\rho), \qquad
  \vect{E} = -\pderiv{\vect{A}}{t} = -\dot\kappa\,a_\phi(\rho)\,\uvec{\phi},
  \quad\text{i.e.}\quad
  E_x = \dot\kappa\,a_\phi\,\frac{y - y_\mathrm{c}}{\rho}, \quad
  E_y = -\dot\kappa\,a_\phi\,\frac{x - x_\mathrm{c}}{\rho}.
\end{equation}
They satisfy Faraday's law: $\nabla\times\vect{E} = -\dot\kappa\,\nabla\times(a_\phi\uvec{\phi})
= -\dot\kappa\,b\,\uvec{z} = -\partial\vect{B}/\partial t$, by the second check above. The
induced field is constant in time and circles the coil's axis; the magnetic field grows in
proportion to $t$.

\subsection{The motion and its conserved quantity}\label{sec:motion}

The ion moves by Newton's equation with $\vect{F} = q(\vect{E} + \vect{v}\times\vect{B})$, in the
plane $F_x = q(E_x + v_yB_z)$, $F_y = q(E_y - v_xB_z)$. Its energy is not conserved: the induced
field does the work $\int q\,\vect{E}\cdot\vect{v}\,\dd t$ (the magnetic force none), carried by
the march as in \xeqref{rf_kick}{eq:work}. The Lagrangian
$L = \tfrac12 m v^2 + q\,\vect{v}\cdot\vect{A}(\vect{x}, t)$ (with no scalar potential: the fields
\eqref{eq:fields} are $-\partial\vect{A}/\partial t$ and $\nabla\times\vect{A}$) is unchanged by a
rotation about the coil's axis, because $\vect{A}$ is $A_\phi(\rho)\,\uvec{\phi}$ at every
moment. So the canonical angular momentum about the axis is conserved:
\begin{equation}\label{eq:canonical}
  J = (x - x_\mathrm{c})\,p_y - (y - y_\mathrm{c})\,p_x + q\,\rho\,\kappa(t)\,a_\phi(\rho)
  = \text{const} = \val{jZero},
\end{equation}
its value at the launch, when $\kappa = 0$. It is exact, so it checks the march at every step.

\subsection{The march}\label{sec:march}

There is no closed form. Newton's equation, with the work as one more state, is marched by
\xeqref{foundations}{eq:rk4} at $h = \val{H}$, the fields \eqref{eq:fields} evaluated at every
stage. The events (the detector, the wire, the world's edge) are located on the cubic Hermite
interpolant by bisection and the error estimated against the step $2h$, as in
\xsecref{first_coil}{sec:method}.

\subsection{The bend to first order in the rate}\label{sec:first}

Along the undeflected line $x = x_0 + v_0t$, $y = y_0$ the transverse force is proportional to
the rate:
\begin{equation}\label{eq:first-order}
  F_y = q\dot\kappa\left(-a_\phi(\rho)\,\frac{x - x_\mathrm{c}}{\rho} - v_0\,t\,b(\rho)\right),
  \qquad
  y^{(1)}(t) = y_0 + \frac{1}{m}\int_0^t\!\!\int_0^{t'} F_y\,\dd t''\,\dd t',
\end{equation}
so to first order the ion's height at the detector's face is $y_0 + \dot\kappa\,S$ with a
sensitivity $S$ that does not depend on the rate: the faster the ramp, the larger the bend, and
the sign of the rate sets its direction. Integrating \eqref{eq:first-order} (the same
Runge--Kutta method) to $t^* = (x_\mathrm{det} - x_0)/v_0 = \val{tStar}$ gives
$S = \val{dydr}$; the detector's middle, $y = \val{yMid}$, then asks for
$\dot\kappa \approx \val{rateFirst}$. That is a scale, not the answer: the bent ion passes nearer
the coil, where $\abs{b}$ is larger (the table above), so the true bend is stronger and a slower
ramp suffices (\S\ref{sec:others}).

\section{The reference solution by hand}\label{sec:hand}

\begin{procedure}
  \item\label{st:fields} \emph{The fields.} At each point, $\rho$ from the coil's centre,
    $m = 4a\rho/(a + \rho)^2$, $K$ and $E$ by \xeqref{first_coil}{eq:agm}, $b$ by
    \xeqref{first_coil}{eq:loop-field} and $a_\phi$ by \eqref{eq:vector-potential}; the fields by
    \eqref{eq:fields} with $\dot\kappa = \val{rate}$.
  \item\label{st:march} \emph{The march.} \S\ref{sec:march} from the launch; every 4 time units:

    \smallskip
    \centerline{\val{march}}
    \smallskip

    The ion enters the detector through its \val{faceRef} face at $(\val{xRef}, \val{yRef})$ at
    $t = \val{tRef}$, with $\vect{p} = (\val{pxRef}, \val{pyRef})$ (step estimate $\val{errRef}$).
    The coil's strength is then $\kappa = \val{kappaEnd}$.
  \item\label{st:checks} \emph{The checks.} $J$ \eqref{eq:canonical} stays at $\val{jZero}$ to
    $\val{jDev}$ along the march; at the entry it is the mechanical part $\val{jMech}$ plus the
    field's $\val{jField}$. The induced field did the work $\val{workRef}$, equal to the change
    of the kinetic energy to $\val{wDev}$: above the coil it points against the ion's motion
    for this sign of the rate, and the ion arrives at speed $\val{speedEnd}$.
\end{procedure}

\begin{result}[The reference flight]
With $\dot\kappa = \val{rate}$ the ion enters the detector at $(\val{xRef}, \val{yRef})$ at
$t = \val{tRef}$. The level is solved.
\end{result}

\section{Comparison with the game}\label{sec:comparison}

\comparisontable

All trajectory differences are at most $\val{largest}$, the game's verified accuracy (Foundations
§6); our step estimates are below $10^{-12}$. The tolerances are ten times the larger of the two
computations' estimates, or the game's own preview--verification gap where that is larger, set
before comparing. The largest difference is in $x_\mathrm{e}$ for the flight that enters through
the detector's top face: there the entry point moves along the face by the error in height
divided by the path's slope, here a few times the height's.

\section{Other solutions}\label{sec:others}

With no supply placed (the table's rows ``no supply'') the coil carries no current and the ion flies
straight, leaving the world's right edge at $t = \val{offT}$. To first order the bend is odd in
the rate (\S\ref{sec:first}): a positive rate bends the ion up, away from the detector below the
launch line, so only negative rates can solve.

The rates that solve form one interval, found by bisection on the rate (the flight marched at
$h = 0.02$, whose error in the path is about $10^{-9}$, far below the change of a unit of rate,
about $10^{-3}$ in height):
\begin{equation}\label{eq:window}
  \val{edgeLo} \;<\; \dot\kappa \;<\; \val{edgeHi},
\end{equation}
$\val{width}$ wide, $\val{widthFrac}$\,\% of the slider's range. At the faster edge the path
grazes the detector's lower left corner (it enters through the \val{faceLo} face at
$y = \val{yEdgeLo}$); at the slower edge it grazes the upper right corner (through the
\val{faceHi} face at $x = \val{xEdgeHi}$). The game's verdicts on rates about 10 inside and
outside each edge, the table's rows labelled by their rates, agree with \eqref{eq:window}.

\section{What the game computes}\label{sec:engine}

The game takes the coil's field from the closed form \xeqref{first_coil}{eq:loop-field} and its
vector potential from \eqref{eq:vector-potential}, with the series for small $m$, both scaled by
the ramp (\physref{2.2}, \emph{Ramped coils and induction}: the quasi-static model,
\ref{as:quasistatic}); its test M10 is the conservation of \eqref{eq:canonical}. It integrates
the equations of motion with its adaptive eighth-order Runge--Kutta method (\physref{3},
\physref{5}, \physref{6}, \physref{7}). The problem is the same as here, and the differences in
the comparison are the two integrations' errors.

\section{Further reading}\label{sec:reading}

The loop's vector potential: \Jackson{§5.5}. The canonical momentum of a charge in a field:
\Jackson{§12.1}. The steady coil and its field: \levelref{first_coil}; a time-dependent source
before this one: \levelref{rf_kick}.

\levelfileappendix

\end{document}
